\documentclass[11pt,a4paper]{article}

\usepackage{titling}
\usepackage[hidelinks]{hyperref}
\usepackage{graphicx}
\usepackage{grffile}
\usepackage{float}
\usepackage{geometry}
\usepackage{listings}

\newcommand{\subtitle}[1]{
  \posttitle{
    \par\end{center}
    \begin{center}\large#1\end{center}
    \vskip0.5em}
}

\begin{document}


\title{Hyper Perform\\ User Manual}
\subtitle{ Repository: \url{https://github.com/claudioMDS/hyper-perform}}
\begin{figure}
			\centering
			\includegraphics[height=200px]{../Images/CodusMaximus_logo.jpg}
\end{figure}

	
\author{
	\textbf{Developers:} \\
	Claudio Da Silva	\emph{14205892}	\\
	Rohan Chhipa		\emph{14188377}	\\
	Avinash Singh		\emph{14043778}	\\
	Jason Gordon		\emph{14405025}	\\\\
}

\date{\textbf{Updated \today}}

\maketitle
\thispagestyle{empty}
\pagebreak

\tableofcontents
\pagebreak

\section{System Overview}
Many different tools are available for measuring the quality of products made, but very few tools exist which assess the quality of the people making said products. People play a huge role in a project, and trying to monitor each and every one becomes a tedious task which diverts man power away from other more critical tasks. Whether it be for an end of year evaluation, or attempting to assess the current status of a project, generating a report on a staff member can help keep up productivity, as well as get them any help they need in order to resume quality performance. By ensuring that there is constant quality performance from each individual on a project, one can increase project quality as well as reduce project risks such as loss of an important team member during a critical stage of a project's life-cycle. 

\section{System Configuration}
This guide has been made for users using a Linux based operating system. To install the HyperPerform system on the machine you will be required to have an active connection to the internet. Please note that high amounts of data might be consumed. 

\subsection{Docker installation}
To install system with ease and avoid all configurations you can download Docker. Docker can be found at \url{www.docker.com} where guides are made available for installing docker on a particular operating system. If you intend to use Docker to install the HyperPerform system then please ensure you install Docker on your machine.

\subsection{Manual Installation}
The manual installation requires you to download the source code from the GitHub repository. The newest release is highly recommended. To carry out a manual installation please ensure you have Maven and the WildFly application server on your machine. \\ \\
Maven can be downloaded from: \url{maven.apache.org} \\ \\
WildFly can be downloaded from: \url{wildfly.org} \\ \\
Please ensure you download WildFly 10. The HyperPerform system was fully tested on this version of WildFly. Any other version might produce unexpected behaviour. \\ 

\section{Installation}

\subsection{Docker Installation}

Assuming you have docker installed on your machine, simply run the following command in terminal: 

\begin{lstlisting}[language=bash]
docker run hyperperform/HyperPerform 
\end{lstlisting}
This will download the HyperPerform Docker image from DockerHub and run it on you machine. \\ \\
The front end component does not have a Docker image at this point in time. To install the front end component please look at section 3.2.5 for the manual installation.

\subsection{Manual Installation}
This installation guide assumes a Linux Server running Ubuntu 14 or higher:

\subsubsection{WildFly}
Once you have downloaded the WildFly application server please carry out the WildFly installations and add a user. Once this is done proceed to installing PostgreSQL.

%Elevate to sudo permissions:

%\begin{lstlisting}
%sudo -s
%\end{lstlisting}
%Install Java JDK 8:
%\begin{lstlisting}
%aptitude update
%aptitude install --with-recommends software-properties-common
%add-apt-repository ppa:webupd8team/java
%aptitude update
%aptitude --with-recommends install oracle-java8-installer vim
%\end{lstlisting}
%Create a WildFly user:
%\begin{lstlisting}
%adduser --no-create-home --disabled-password --disabled-login wildfly
%\end{lstlisting}
%Download and extract the WildFly installation:
%\begin{lstlisting}
%cd /opt
%wget --tries=0 --continue http://download.jboss.org/wildfly/10.0.0.Final/wildfly-10.0.0.Final.tar.gz
%tar -xzvf wildfly-10.0.0.Final.tar.gz
%ln -s wildfly-10.0.0.Final wildfly
%chown -R wildfly wildfly
%\end{lstlisting}
%\pagebreak
%Copy the initialization scripts to the required folders:
%\begin{lstlisting}
%cp /opt/wildfly/docs/contrib/scripts/init.d/wildfly-init-debian.sh /etc/init.d/wildfly
%update-rc.d /etc/init.d/wildfly defaults
%cp /opt/wildfly/docs/contrib/scripts/init.d/wildfly.conf /etc/default/wildfly
%cd /etc/default
%\end{lstlisting}
%Now to edit the files:
%\begin{lstlisting}
%nano wildfly
%\end{lstlisting}
%Uncomment and/or Edit the following lines:
%\begin{lstlisting}
%JBOSS_HOME="/opt/wildfly"
%JBOSS_USER=wildfly
%JBOSS_MODE=standalone
%JBOSS_CONFIG=standalone-full.xml
%JBOSS_CONSOLE_LOG="/var/log/wildfly/console.log"
%\end{lstlisting}
%Change all 127.0.0.1 in standalone-full.xml to 0.0.0.0, then run the following commands:
%\begin{lstlisting}
%service wildfly start
%cd /opt/wildfly/bin
%./add-user.sh
%\end{lstlisting}

\subsubsection{PostgreSQL}
The install PostgreSQL on your machine: \\\\
Install via terminal:
\begin{lstlisting}
sudo apt-get update
sudo apt-get install postgresql postgresql-contrib
\end{lstlisting}
Then create the required database:
\begin{lstlisting}
sudo -i -u postgres
psql
createdb hyperperform
\end{lstlisting}
To configure PostgreSQL to connect remotely:
\begin{lstlisting}
sudo nano /etc/postgresql/9.3/main/postgresql.conf
\end{lstlisting}
Edit the following lines:
\begin{lstlisting}
listen_addresses = "*"
\end{lstlisting}



\subsubsection{ActiveMQ}
To setup ActiveMQ on your server: \\
\begin{itemize}
	\item Start up your WildFly application server
	\item Navigate to WildFly management console on localhost:9990
	\item Navigate to configurations tab and click on sub-systems
	\item Scroll down and search for Messaging-ActiveMQ and click on it
	\item Click on default, select queues/topics
	\item Click add and input the following information:
		\begin{itemize}
			\item Name*: hyperperform
			\item JNDI Names*: java:/jms/queue/hyperperform
		\end{itemize}
	\item Click save \\
\end{itemize}

\subsubsection{Deploying to WildFly}
To deploy the HyperPerform system to the application you will need to build the system using from the source code. 

\begin{itemize}
	\item Ensure the WildFly server is running.
	\item Navigate to \url{https://github.com/HyperPerform/hyper-perform-server/releases} and download the newest release source code.
	\item Extract the source code
	\item Navigate to the root directory of the source code. A file named pom.xml should be clearly visible.
	\item Run the following command: mvn clean wildfly:deploy
	\item Maven will automatically deploy the compiled code to WildFly
\end{itemize}



\subsubsection{Front-end Dashboard}
To start up the front end please ensure you have Node installed on your machine. Node can be found at \url{https://nodejs.org/en/}. \\

\begin{itemize}
	\item Download the Dashboard source code from \url{https://github.com/HyperPerform/hyper-perform-web-application}
	\item Navigate to the root directory of the source code
	\item Run the following command in terminal: npm install
	\item Next run: npm install gulp
	\item Finally: gulp serve
\end{itemize}

%\pagebreak

%\section{Installation of HyperPerform}
%
%\subsection{Server}
%To setup the back end, take the hyperperform.war and deploy it on the WildFly management console
%
%\subsection{Dashboard}
%To install the dashboard on your server:
%\begin{lstlisting}
%npm install hyperperform
%\end{lstlisting}


\end{document}